\section{Theoretical Analysis}
\label{sec:theory}

The empirical picture of Section~\ref{sec:results} is consistent and, in two
places, exact: the advantage is positive and interval-separated at every
capacity, it is largest where capacity is scarce, and it is precisely zero on
the homogeneous control. This section asks why those three facts hold as a
matter of the scorer's algebra rather than as an empirical accident, and why the
benefit, though real, is bounded in a way that is congruent with the partial H1
verdict. We derive closed forms where the multiplicative structure of
Eq.~\ref{eq:capg} permits one and fall back to a qualitative mechanism where it
does not. No quantity stated here is new: every numeric anchor is the frozen
value already reported in Section~\ref{sec:results}.

\subsection{When the multiplier can reorder a pair}
\label{sec:theory_reorder}

Fix a seed and consider two candidate pairs $(h,c)$ and $(h',c')$ with base
scores $S\equiv S_{\mathrm{base}}(h,c)$ and $S'\equiv S_{\mathrm{base}}(h',c')$
and host context scores $a\equiv\mathrm{ACS}(h)$ and $a'\equiv\mathrm{ACS}(h')$.
Under CAP-G the pair $(h,c)$ outranks $(h',c')$ iff
%
\begin{equation}
\label{eq:reorder_condition}
S\,\bigl(1+\rho a\bigr)\;>\;S'\,\bigl(1+\rho a'\bigr).
\end{equation}
%
Suppose the base scorer ranks them the other way, $S'>S>0$, so that CAP-G
reorders the pair exactly when Eq.~\ref{eq:reorder_condition} holds. Dividing
through by $S'(1+\rho a')>0$ and rearranging gives the clean reordering
criterion
%
\begin{equation}
\label{eq:reorder_ratio}
\frac{S}{S'}\;>\;\frac{1+\rho a'}{1+\rho a}.
\end{equation}
%
The left side is the base-score ratio, which lies in $(0,1)$ by assumption; the
right side is below $1$ precisely when $a>a'$, that is, when the lower-base pair
sits on the higher-context host. So CAP-G promotes a pair over a higher-scored
competitor only if it has both higher context and a base-score ratio close
enough to one to be overturned by the context gap. Writing the base scores as
close, $S'=S(1+\epsilon)$ with small $\epsilon>0$, a first-order expansion of
Eq.~\ref{eq:reorder_ratio} yields the reordering condition
%
\begin{equation}
\label{eq:reorder_firstorder}
\rho\,(a-a')\;>\;\epsilon\;+\;O(\epsilon^2,\rho^2),
\end{equation}
%
which states the mechanism compactly: \emph{a reordering occurs only among pairs
whose base scores are separated by less than the context-weighted gap
$\rho(a-a')$}. Pairs whose base scores are far apart, $\epsilon\gg\rho(a-a')$,
keep their order regardless of context. The context factor therefore acts only
in the neighborhood of base-score ties, and it acts there in the direction of
the more mission-critical host. This is the algebraic content of the H4 result
(Section~\ref{sec:results}): because CAP-G perturbs the order only within
near-tied bands rather than rewriting it wholesale, its context-blind precision
is statistically indistinguishable from the base scorer's (difference
$+0.1$~pp, interval $[-1.5,+1.6]$), since the wholesale ordering on the blind
target is left intact.

\subsection{Why heterogeneity is the lever}
\label{sec:theory_heterogeneity}

Eq.~\ref{eq:reorder_firstorder} makes the benefit a function of context
\emph{gaps} $a-a'$ between near-tied pairs, not of context \emph{levels}. The
number of order-changing pairs, and hence the mission-precision the reordering
can recover, grows with how widely the context score is spread across the fleet.
Let the fleet context scores have mean $\bar a$ and variance
$\sigma_a^2=\operatorname{Var}_h[\mathrm{ACS}(h)]$. For a randomly drawn pair of
near-tied hosts the expected squared context gap is
%
\begin{equation}
\label{eq:gap_variance}
\mathbb{E}\bigl[(a-a')^2\bigr]\;=\;2\,\sigma_a^2,
\end{equation}
%
so the typical magnitude of the reordering term $\rho(a-a')$ in
Eq.~\ref{eq:reorder_firstorder} scales with the context standard deviation
$\sigma_a$. As $\sigma_a$ rises, more near-tied bands contain a context gap large
enough to clear the local base-score separation $\epsilon$, more mission-critical
pairs are lifted into the top-$K$, and the mission-weighted precision gain over
the base scorer grows. The gain is thus governed by the heterogeneity (variance)
of the asset context score across the fleet: it is monotone increasing in
$\sigma_a$ from a floor of zero, with the exact scaling of the realized
precision lift through that variance depending on the local density of base
scores near each tie, which is fleet-specific and not available in closed form.
We state that residual dependence qualitatively rather than assign it a constant
the data do not pin down. The direction, however, is unambiguous and is exactly
what the mechanism control measures: dispersion in context is the resource CAP-G
converts into mission precision.

\subsection{The homogeneous limit: an exact zero}
\label{sec:theory_homogeneous}

The variance argument has a sharp boundary case. On a homogeneous fleet every
host carries the same context score, $\mathrm{ACS}(h)\equiv a_0$, so
$\sigma_a^2=0$ and every context gap $a-a'$ in
Eq.~\ref{eq:reorder_firstorder} is identically zero. The CAP-G multiplier
collapses to a single global constant,
%
\begin{equation}
\label{eq:homogeneous}
\begin{split}
S_{\mathrm{CAP\text{-}G}}(h,c) &= S_{\mathrm{base}}(h,c)\,\bigl(1+\rho a_0\bigr) \\
&= \kappa\,S_{\mathrm{base}}(h,c),\qquad \kappa=1+\rho a_0>0,
\end{split}
\end{equation}
%
and multiplying every pair's score by the same positive $\kappa$ is an
order-preserving transformation: it cannot change which $K$ pairs are ranked
highest, so every rank-based metric, MWP@K included, is identical to the base
scorer's at every capacity. The gain is therefore not merely small but
\emph{exactly zero}, for every $K$, as an algebraic identity rather than an
empirical near-miss. This is precisely the H3 mechanism-control result of
Section~\ref{sec:results}, where the realized homogeneous-fleet advantage is
$0.0$~pp at every $K$ (interval $[0,0]$); the derivation here shows that outcome
was forced by the multiplicative form of Eq.~\ref{eq:capg} and not contingent on
the choice of constants. Conversely, any non-zero advantage observed on the
heterogeneous fleet must originate in $\sigma_a>0$, which is the sense in which
the heterogeneous gain is attributable to fleet dispersion and not to the
scoring transform.

\subsection{Why the benefit is bounded}
\label{sec:theory_bounded}

The same algebra that makes the gain positive also caps it, and the cap is
consistent with reporting H1 as partial. Three bounds compound. First, by
Eq.~\ref{eq:reorder_firstorder} the multiplier touches only the near-tied bands;
pairs whose base scores are well separated are immovable, so the reachable set of
reorderings is a strict subset of all pairs and the precision the reordering can
recover is bounded above by the mission mass sitting inside those bands. Second,
the context factor is itself bounded: $\mathrm{ACS}(h)\in[0,1]$ gives
$1+\rho\,\mathrm{ACS}(h)\in[1,\,1+\rho]$, so the largest relative re-weighting any
pair can receive is $1+\rho$, and the most a low-base mission-critical pair can be
promoted over a high-base competitor is fixed by that finite ratio rather than
being unbounded. Third, the benefit interacts with capacity. As $K$ grows the
top-$K$ set absorbs most high-base pairs regardless of order, so the marginal
mission mass that reordering can still rescue shrinks; this is the capacity-decay
mechanism whose monotone realization is the H2 result, and it drives the observed
decline of the advantage from $+9.5$~pp at $K=50$ to $+3.0$~pp at $K=100$ to
$+0.7$~pp at $K=250$.

These three bounds together explain the shape of the primary finding without
softening it. The advantage is genuinely positive and interval-separated at every
capacity, and it is largest exactly where remediation budget is scarcest, but
because it decays toward zero as capacity grows, its average across the three
pre-registered capacities is a bounded quantity, and the realized average of
$4.4$~pp sits just below the locked $5$~pp bar. The theory predicts a positive but
capacity-attenuated effect that is strongest at small $K$; that is exactly the
partial verdict the frozen results record, and we present it as a bounded effect
rather than a stronger claim than the data support. Framed in the information
terms of~\cite{cover2006elements}, the context dimension can only reorder pairs to
the extent that mission value carries information the base ranking has not already
expressed; where the budget is large enough to cover the high-base pairs anyway,
that incremental information has little left to act on, and the bounded gain is the
visible signature of a bounded information increment.

\subsection{Robustness and Sensitivity}
\label{sec:theory_robustness}

\textbf{An adversary who hides in low-context hosts.} The reordering analysis
exposes a specific worst case for a context-aware queue. Section~\ref{sec:results}
showed the dual of this for the base scorer's hygiene channel; here we reason about
the mission channel. Suppose a strategic attacker, knowing the defender ranks with
CAP-G, deliberately concentrates exploitation on \emph{low}-context hosts, where
$\mathrm{ACS}(h)$ is small and the multiplier $1+\rho\,\mathrm{ACS}(h)$ is nearest
its floor of $1$. On exactly those hosts the context factor gives the least lift,
so a compromise there is the case CAP-G is least able to surface relative to the
base scorer. Two structural facts bound the damage. First, the mission-weighted
objective is, by design, indifferent to such hosts in proportion to their low
mission value: a successful compromise of a low-context host costs little
mission-weighted exposure precisely because the MCTP target counts a host as
mission-critical only when its ACS exceeds the fleet 75th percentile
(Section~\ref{sec:eval}), so the attack lands where the loss the objective is
defined to minimize is smallest. The attacker buys access to the hosts the
multiplier ignores by accepting the hosts the mission objective also discounts;
the two are the same hosts by construction. Second, on those low-context pairs CAP-G
reduces toward the base scorer, since when $\mathrm{ACS}(h)\!\to\!0$ the factor
$1+\rho\,\mathrm{ACS}(h)\!\to\!1$ and Eq.~\ref{eq:capg} returns
$S_{\mathrm{base}}$; the defender therefore retains the full hygiene-augmented base
ranking on exactly the hosts where context adds nothing, rather than being left
defenseless. The residual risk is real and we do not minimize it: an attacker
willing to operate only against genuinely low-value hosts faces a queue that has
correctly deprioritized them, which is the intended behavior of a mission-weighted
objective, but a defender who disputes the mission valuation, for instance one who
believes a nominally low-context host is a stepping stone to a critical one, should
revisit the context assignment itself rather than the scorer, because CAP-G faithfully
propagates whatever mission values the fleet model encodes.

\textbf{Sensitivity to the context emphasis.} The bounds above are stated for
fixed $\rho$, and the realized study used the locked-grid value $\rho{=}8.0$
selected on five held-out calibration seeds with the pre-registered stop rule
(Section~\ref{sec:method}). Because the multiplier range $1+\rho$ is monotone in
$\rho$, larger $\rho$ widens the band of base-score ties that context can overturn
in Eq.~\ref{eq:reorder_firstorder}, which is consistent with the monotone rise of
calibration MWP@50 across the locked grid reported in Section~\ref{sec:method};
the chosen boundary value is therefore a conservative lower bound on context
emphasis, and the bounded-gain conclusion is if anything understated with respect
to $\rho$. Sensitivity to \emph{which} context dimension drives the effect is
settled empirically by the ablation of Section~\ref{sec:results}, where removing
asset criticality is the only deletion whose MWP@50 change separates from zero;
the variance lever $\sigma_a$ of Section~\ref{sec:theory_heterogeneity} is, on this
target, carried predominantly by the criticality dimension, and the network-zone
and data-sensitivity terms contribute little dispersion that the mission target
rewards. We change none of those calibration or ablation values here; the theory
explains their direction and leaves their magnitudes to the frozen results.
